\begin{idstatement}{D.2.5}{Proved unrestricted sign under trivial monodromy}
\label{st:D.2.5}
Let $f:Y\to\PP^1$ be a proper surjective holomorphic map with
connected fibres from a compact connected complex threefold with
$b_1(Y)=b_2(Y)=0$. Suppose its smooth general fibre is a complex
two-torus and its monodromy representation on $H_1(F,\Q)$ is
trivial. Then
\[
 c_3(Y)\le0,\qquad b_3(Y)\ge2.
 \tag{D.S.1}
\]
The fibres may have arbitrary singularities and arbitrary
multiplicities. In particular this proves the sign assertion of
D.2 for a holomorphic algebraic reduction; no existence assertion
is required for the conclusion.
\end{idstatement}
\begin{proof}
Assume $c_3(Y)>0$. Poincar\'e duality gives
$c_3(Y)=2-b_3(Y)$, and the alternating nondegenerate intersection
pairing on $H^3(Y,\Q)$ makes $b_3(Y)$ even. Hence $b_3(Y)=0$.
Apply C4c to the Picard component
$P=K_Y\otimes\Pic^0(Y)$. It gives a fixed critical full fibre
$F$ such that
\[
 H^0\bigl(F,\Omega_Y^1|_F\otimes M|_F\bigr)\ne0
 \qquad(M\in\Pic^0(Y)).
 \tag{D.S.2}
\]
Indeed the Serre-dual parameter $K_Y\otimes L^{-1}$ runs through
$\Pic^0(Y)$ when $L$ runs through $P$.

Statement A.3.5 excludes reducible support in (D.S.2), including
all unequal-multiplicity possibilities. Thus $F=mS$ for one
irreducible reduced Cartier divisor $S$ and one positive integer
$m$. We specify why a smooth $S$ is excluded under the present
\emph{rational} Betti-number assumptions. D.1.9 supplies
$\beta\in H^1(Y,\OO_Y)$ whose restriction near every full fibre
comes from a nonzero integral transgression class. The line
$M_r=\exp(2\pi i r\beta)$, for real irrational $r$, has that
unitary local realization. The normalized root change $t=s^m$
near a smooth $S$ is the finite \emph{\'etale} map of A.3.4:
in a chart $t=uz^m$ it has equations $w^m=u$, $s=wz$.
Its central reduction $S'$ is smooth, and its nearby fibres are
the original tori. Proper smooth local triviality transports
the vanishing of degree-zero and degree-one cohomology of the
pulled-back unitary local system from a nearby torus to $S'$.
The injectivity A.4.7 and the conormal sequence then give zero
ambient cotangent sections on $S'$. Each layer of $mS'$ has the
same twist, since $s$ trivializes its conormal. The finite
filtration and exact invariants functor for the cyclic cover
give $H^0(mS,\Omega_Y^1|_{mS}\otimes M_r)=0$, contradicting
(D.S.2). This proof uses only the rational version of D.1.9.
Therefore $S$ is singular. A.5.2, applied using
$\omega_S=(K_Y\otimes\OO_Y(S))|_S$ and
$\chi(\OO_S)=0$, makes $S$ nonnormal.

The ideal filtration of $mS$ has the finitely many locally free
graded factors on $S$
\[
 E_j=(\Omega_Y^1\otimes\OO_Y(-jS))|_S,
 \qquad 0\le j<m.
 \tag{D.S.3}
\]
If all the groups $H^0(S,E_j\otimes M|_S)$ were zero, the exact
filtration sequences would make (D.S.2) zero. Thus their
nonvanishing loci cover $\Pic^0(Y)$. Each locus is closed
complex analytic by the actual Picard family and proper
semicontinuity. The group $\Pic^0(Y)$ is a connected complex
vector space, since $H^1(Y,\Z)=0$. A finite union of proper
closed analytic subsets cannot cover it. Some fixed $j$
therefore satisfies
\[
 H^0(S,E_j\otimes M|_S)\ne0
 \qquad(M\in\Pic^0(Y)).
 \tag{D.S.4}
\]

Let $\eta:T\to S$ be the finite normalization and
$\pi:Z\to T$ its minimal resolution, and put $h=\eta\pi$,
viewed also as a holomorphic map to $Y$.
Pullback of sections in (D.S.4) is injective: $S$ is reduced
and irreducible, and a section of the locally free sheaf $E_j$
vanishing on its dense smooth locus is zero. Thus the fixed
vector bundle $h^*(\Omega_Y^1\otimes\OO_Y(-jS))$ has a
section after every twist by $h^*M$, $M\in\Pic^0(Y)$.
For any Gauduchon metric $\omega$ on $Z$, the uniform
negative-degree bound A.5.6, applied to all positive and
negative powers of $M$, gives
\[
 \deg_\omega h^*M=0\qquad(M\in\Pic^0(Y)).
\]
Every line bundle $G$ on $Y$ has a positive power in
$\Pic^0(Y)$, because $H^2(Y,\Z)$ is finite. Additivity of
degree therefore gives $\deg_\omega h^*G=0$ for every such
$G$. No classification of the surface $Z$ is used in this step.

The conductor comparison A.5.3 gives an effective divisor
$D_Z$ with
\[
 \omega_Z(D_Z)\simeq h^*(K_Y\otimes\OO_Y(S)).
 \tag{D.S.5}
\]
Moreover $D_Z\ne0$. Indeed a reduced hypersurface is
Cohen--Macaulay and hence satisfies Serre's condition $S_2$;
its nonnormality must therefore occur in codimension one.
The divisorial conductor is nonzero, and its strict transform
is part of $D_Z$.
Taking Gauduchon degree in (D.S.5) yields
\[
 \deg_\omega K_Z=-\deg_\omega\OO_Z(D_Z)<0.
\]
A nonzero holomorphic section of $K_Z$ would have an effective
zero divisor, whose degree is nonnegative. This contradiction
proves $p_g(Z)=h^0(Z,K_Z)=0$.

Choose a small closed fibre neighbourhood $U=f^{-1}(\Delta)$.
Embedded resolution of the divisor $F$ gives a proper
holomorphic modification
\[
 q:\widehat U\longrightarrow U
\]
which is a finite composite of blowups with smooth centres in
the successive central divisors, is an isomorphism over
$\Delta^*$, has smooth total space, and makes the reduced
support of $(f\circ q)^{-1}(0)$ a simple normal-crossings
divisor. We use the analytic embedded-resolution theorem in
\cite[Theorem~1.6]{BierstoneMilman1997}; properness of the
central fibre makes its locally finite construction finite
after shrinking $\Delta$. The strict transform of $S$ is a
resolution of its normalization and has geometric genus zero
by birational invariance of holomorphic two-forms on smooth
compact surfaces. A new exceptional component is $\PP^2$ over
a point or a $\PP^1$-bundle over a smooth compact curve.
Both have geometric genus zero: restriction of a holomorphic
two-form to a ruling chart has negative fibre degree, and
$H^0(\PP^2,K_{\PP^2})=0$. Further strict transforms preserve
this property. Consequently every component of the final
central reduced divisor has geometric genus zero.

The map $f\circ q$ remains proper with connected fibres:
a blowup has connected fibres and its structure-sheaf direct
image is the original structure sheaf. Its smooth fibres over
$\Delta^*$ are the original two-tori under the isomorphism
$q|_{\Delta^*}$, so their $H^1$ monodromy remains the identity.
The multiplicities of the resolved central fibre are arbitrary
positive integers and are all permitted in D.2.3. That
proposition now gives nontrivial monodromy on $H^2$ and hence
on $H^1$, a contradiction. This proves (D.S.1).
\end{proof}

The sign result uses the analytic logarithmic comparison of
Steenbrink. The explicit compact construction in D.2.7 proves
that the stated class is nonempty and attains $c_3=0$.
Together these results settle both the existence and the sign
questions in D.2.


% D.2.6 from part_D_periods.tex

